%%% FIELD proof-cell-gaussian
Put \(s_{ja}=p_je_{ja}\) and \(s_0=p_{\min}e_{\min}>0\).  Consistency and conditional randomization give
\[
 E_{P_n}(Y\mid X=j,A=a)=E_{P_n}\{Y(a)\mid X=j\}=\mu_{n,ja}.
\]
Thus, with \(I_{i,ja}=\mathbf 1\{X_{ni}=j,A_{ni}=a\}\), \(U_{i,ja}=I_{i,ja}(Y_{ni}-\mu_{n,ja})\), \(\widehat s_{ja}=n^{-1}\sum_iI_{i,ja}\), and \(\overline U_{ja}=n^{-1}\sum_iU_{i,ja}\), one has \(E U_{i,ja}=0\), \(|U_{i,ja}|\le2B_Y\), and, on \(\widehat s_{ja}>0\),
\[
 \sqrt n(\widehat\mu_{n,ja}-\mu_{n,ja})
 =n^{-1/2}\sum_{i=1}^n\frac{U_{i,ja}}{s_{ja}}
 +R_{n,ja},\qquad
 R_{n,ja}=\sqrt n\,\overline U_{ja}
       \left(\frac1{\widehat s_{ja}}-\frac1{s_{ja}}\right).
\]
Let \(\mathcal A_n=\bigcap_{j,a}\{|\widehat s_{ja}-s_{ja}|\le s_{ja}/2\}\).  Hoeffding's inequality, obtained here by applying the exponential Markov bound to the Bernoulli indicators and optimizing its scalar parameter, gives
\[
 P_n(\mathcal A_n^c)\le2JK\exp(-ns_0^2/2).
\]
On \(\mathcal A_n\), Cauchy--Schwarz and the variance identities for the two sample averages yield
\[
 E\{|R_{n,ja}|\mathbf1_{\mathcal A_n}\}
 \le \frac{2\sqrt n}{s_{ja}^2}
       \{E\overline U_{ja}^2\}^{1/2}
       \{E(\widehat s_{ja}-s_{ja})^2\}^{1/2}
 \le \frac{2B_Y}{s_0^2\sqrt n}.
\]
The convention on an empty cell changes a bounded-Lipschitz expectation by at most its probability.  Since \(\sup_{x>0}\sqrt{x}\exp(-s_0^2x/2)=(e^{-1/2}/s_0)\), stacking all cells and taking action-zero differences therefore gives
\[
 d_{\mathrm{BL}}\!\left\{\mathcal L(E_n),
 \mathcal L\!\left(n^{-1/2}\sum_{i=1}^n\xi_n(O_{ni})\right)\right\}
 \le \frac{C_{\mathrm{rat}}}{\sqrt n},
 \quad
 C_{\mathrm{rat}}=8JK\sqrt d\,B_Ys_0^{-2}+2JK(e^{-1/2}/s_0).
\]

For completeness, the finite-dimensional normal comparison needed next can be derived directly.  If independent centered vectors \(X_i\in\mathbb R^d\) satisfy \(\sum_iE(X_iX_i^\top)=\Sigma\), \(\lambda_{\min}(\Sigma)\ge v_{min}\), and \(\lambda_{\max}(\Sigma)\le v_{max}\), Gaussian interpolation of a bounded one-Lipschitz test function, followed by Taylor's formula with integral remainder in each leave-one-out vector, gives
\[
 d_{\mathrm{BL}}\{\mathcal L(\textstyle\sum_iX_i),N(0,\Sigma)\}
 \le 32(d+1)^{3/2}\max(1,\sqrt{v_{\max}}),v_{\min}^{-3/2}
       \sum_iE\|X_i\|_2^3.                                           \tag{1}
\]
Here is the derivative calculation behind (1).  After replacing \(X_i\) by \(\Sigma^{-1/2}X_i\), solve the Gaussian interpolation equation by
\[
 g_f(x)=-\int_0^\infty
 \left[E f\{e^{-s}x+(1-e^{-2s})^{1/2}Z\}-Ef(Z)\right]ds,qquad Z\sim N(0,I_d).
\]
Gaussian integration by parts changes \(E\{\Delta g_f(W)-W^\top\nabla g_f(W)\}\) into the sum of the leave-one-out second-order Taylor remainders.  Splitting the integral at \(s=\|X_i\|_2^2\), using \(|f(x)-f(y)|\le\|x-y\|_2\) before the split and the Gaussian derivative formula after it, bounds the \(i\)-th remainder by \(32(d+1)^{3/2}E\|X_i\|_2^3\).  Undoing the whitening multiplies this by at most \(v_{\min}^{-3/2}\max(1,\sqrt{v_{\max}})\), which proves (1), including for nonsmooth Lipschitz tests by Gaussian mollification and monotone passage to the limit.

In the present problem \(E\xi_n(O)=0\), \(\operatorname{Var}\{\xi_n(O)\}=V_n\), and
\[
 \|\xi_n(O)\|_2\le B_\xi:=\frac{4B_Y\sqrt d}{p_{\min}e_{\min}}.
\]
Apply (1) to \(X_i=\xi_n(O_{ni})/\sqrt n\).  Then \(\sum_iE\|X_i\|_2^3\le B_\xi^3/\sqrt n\).  The triangle inequality gives the asserted result with the explicit choice
\[
 C_{\mathrm{BE}}=C_{\mathrm{rat}}+
 32(d+1)^{3/2}\max(1,\sqrt{v_{\max}})v_{\min}^{-3/2}B_\xi^3.
\]
Every denominator is protected by \(s_0>0\), and the covariance inversion is protected by \(v_{\min}>0\), as required.

%%% FIELD proof-bregman-identity
Write \(q(v)=q_{c,b}(v)\) and \(\Psi(v)=\Psi_{c,b}(v)\).  Positivity of every coordinate of \(b\) makes all logarithms finite, and direct differentiation gives
\[
 \nabla\Psi(v)=q(v),\qquad \nabla^2\Psi(v)=c^{-1}C\{q(v)\}.
\]
Moreover,
\[
 \log\frac{q(u+z)_a}{q(u)_a}
 =\frac{z_a}{c}-\log\frac{\sum_rb_r e^{(u_r+z_r)/c}}
                              {\sum_rb_r e^{u_r/c}}.
\]
Multiplication by \(q(u+z)_a\), summation over \(a\), and multiplication by \(c\) give exactly
\[
 c\,\mathrm{KL}\{q(u+z)\|q(u)\}
 =\Psi(u)-\Psi(u+z)+z^\top q(u+z)=\ell_{c,b}(u,z).
\]
For \(f(s)=\Psi(u+sz)\), twice applying the fundamental theorem of calculus yields
\[
 f(0)-f(1)+f'(1)=\int_0^1s f''(s)\,ds,
\]
and hence
\[
 \ell_{c,b}(u,z)=c^{-1}\int_0^1s\,z^\top C\{q(u+sz)\}z\,ds.
\]
The covariance matrix \(C(r)\) is positive semidefinite for every probability vector \(r\), proving nonnegativity.  Also, convexity and the elementary bounds \(\min_a z_a\le z^\top r\le\max_a z_a\) imply
\[
 \ell_{c,b}(u,z)\le \max_a z_a-\min_a z_a=\operatorname{osc}(z).
\]
Finally, \(z^\top C(r)z=\operatorname{Var}_r(z_a)\le\operatorname{osc}(z)^2/4\); inserting this bound into the integral and using \(\int_0^1s\,ds=1/2\) gives \(c\ell_{c,b}(u,z)\le\operatorname{osc}(z)^2/8\).

%%% FIELD proof-exact-regret-path
Conditional randomization and consistency identify \(\mu_{n,ja}=E(Y\mid X=j,A=a)\); overlap makes each conditioning event positive.  In stratum \(j\), subtracting the action-zero mean, multiplying the objective by \(\sqrt n\), and using \(\lambda_n(c)=c/\sqrt n\) gives
\[
 \sqrt n\left\{\sum_a\mu_{n,ja}\pi_{ja}-\lambda_n(c)\mathrm{KL}(\pi_j\|b_j)-\mu_{n,j0}\right\}
 =(Ru_{n,j})^\top\pi_j-c\,\mathrm{KL}(\pi_j\|b_j).
\]
The Lagrange equations for maximizing the right side over the simplex are
\[
 (Ru_{n,j})_a-c\{1+\log(\pi_{ja}/b_{ja})\}=\kappa_j,
\]
and benchmark positivity gives the unique solution \(q_{c,b_j}(Ru_{n,j})=\pi^*_{n,c,j}\).  For arbitrary \(v,w\in\mathbb R^K\), direct substitution of the Gibbs probabilities shows
\[
 \big[v^\top q(v)-c\mathrm{KL}\{q(v)\|b\}\big]
 -\big[v^\top q(w)-c\mathrm{KL}\{q(w)\|b\}\big]
 =c\mathrm{KL}\{q(w)\|q(v)\}=\ell_{c,b}(v,w-v).
\]
Set \(v=Ru_{n,j}\) and \(w=R\widehat u_{n,j}=R(u_{n,j}+E_{n,j})\), sum with weights \(p_j\), and divide by \(\sqrt n\).  This gives
\[
 \sqrt n\,\mathscr R_n(c)=\sum_jp_j\ell_{c,b_j}(Ru_{n,j},RE_{n,j}).
\]
Multiplication by \(1+c\), followed by \(t=c/(1+c)\), is exactly the interior definition of \(G(u_n,E_n)(t)=T_n(t)\).  The endpoint equalities are compactified limits proved in the hard- and high-endpoint lemmas; they are not additional observed-law definitions of the causal regret.

%%% FIELD proof-high-endpoint
For \(v\in\mathbb R^K\), differentiate the exponential tilt \(s\mapsto q_{c,b}(sv)\):
\[
 \frac{d}{ds}q_{c,b}(sv)=c^{-1}C\{q_{c,b}(sv)\}v.
\]
Since \(\sum_a q_a|v_a-q^\top v|\le2\|v\|_\infty\), integration from \(0\) to \(1\) gives \(\|q_{c,b}(v)-b\|_1\le2\|v\|_\infty/c\).  For probability vectors \(r,b\), expansion of the diagonal and rank-one terms gives
\[
 |z^\top\{C(r)-C(b)\}z|
 \le3\|z\|_\infty^2\|r-b\|_1.
\]
The integral representation of the Bregman loss therefore implies
\[
 \begin{aligned}
 \left|c\ell_{c,b}(u,z)-\tfrac12z^\top C(b)z\right|
 &\le3\|z\|_\infty^2\int_0^1s\|q_{c,b}(u+sz)-b\|_1ds\\
 &\le \frac{3\|z\|_\infty^2(\|u\|_\infty+\|z\|_\infty)}{c}.
 \end{aligned}
\]
Also \(\ell_{c,b}(u,z)\le\operatorname{osc}(z)^2/(8c)\), so replacing \(c\ell\) by \((1+c)\ell\) adds only a further computable \(O(c^{-1})\) term on bounded sets.  This proves the asserted endpoint modulus.

For \(x\in\mathbb R^{K-1}\), the categorical variance identity gives
\[
 x^\top R^\top C(b)Rx
 =\sum_{0\le a<r\le K-1}b_ab_r\{(Rx)_a-(Rx)_r\}^2
 \ge b_0\sum_{a=1}^{K-1}b_ax_a^2
 \ge\beta^2\|x\|_2^2.
\]
Thus the projected endpoint is positive definite under the declared support condition.

%%% FIELD proof-hard-endpoint
Let \(M=\max_av_a\), \(I=I(v)\), and \(B_I=\sum_{a\in I}b_a\).  Benchmark support gives \(B_I\ge\beta\), while
\[
 \Psi_{c,b}(v)=M+c\log\left\{B_I+\sum_{a\notin I}b_ae^{-(M-v_a)/c}\right\}.
\]
Consequently \(|\Psi_{c,b}(v)-M|\le c\log(1/\beta)\).  The total Gibbs mass outside \(I\) is at most \(\beta^{-1}e^{-\gamma(v)/c}\); conditional on \(I\), its proportions are exactly proportional to \(b_a\).  Hence
\[
 \|q_{c,b}(v)-q_{0,b}(v)\|_1
 =2q_{c,b}(v)(I^c)
 \le2\beta^{-1}e^{-\gamma(v)/c}.
\]
Apply the log-sum bound at \(u\) and \(u+z\), and the last display at \(u+z\), in the definitions of \(\ell_c\) and \(\ell_0\).  The triangle inequality gives
\[
 |\ell_{c,b}(u,z)-\ell_{0,b}(u,z)|
 \le2c\log(1/\beta)+2\|z\|_\infty\beta^{-1}e^{-\gamma(u+z)/c}.
\]
This includes every exact multiway tie; when all coordinates tie, \(\gamma=\infty\) and the Gibbs and hard lotteries agree exactly.

For the probability bound, each within-stratum pairwise difference is \(m+v^\top Z_V\) for a nonzero contrast vector \(v\).  If one action is zero then \(\|v\|_2=1\), and otherwise \(\|v\|_2=\sqrt2\).  Since \(V\succeq v_{\min}I_d\), its variance is at least \(v_{\min}\).  A univariate normal variable of variance at least \(v_{\min}\) has density at most \((2\pi v_{\min})^{-1/2}\), so
\[
 P\{|m+v^\top Z_V|\le a\}\le\frac{2a}{\sqrt{2\pi v_{\min}}}.
\]
There are \(J K(K-1)/2\) pairs.  The union bound proves the displayed \(JK(K-1)a/\sqrt{2\pi v_{\min}}\) inequality, uniformly in the means and covariance.

%%% FIELD proof-path-wellposedness
For fixed \((u,z)\), the two endpoint lemmas show that the interior formula for \(G(u,z)\) converges to the declared values at both \(t=0\) and \(t=1\); hence it is in \(C([0,1])\).  Every rational-time evaluation is a Borel function of \((u,z)\), including the hard evaluation because it is a finite combination of comparisons and benchmark-weighted finite sums.  Since \(C([0,1])\) is separable and its Borel sigma-field is generated by rational-time evaluations, \((u,z)\mapsto G(u,z)\) is Borel.

The global loss bounds imply, for every \(t\),
\[
 0\le G(u,z)(t)\le\sum_jp_j\left\{\operatorname{osc}(Rz_j)
       +\operatorname{osc}(Rz_j)^2/8\right\}.
\]
The right side is finite for every finite \(z\), proving almost-sure finiteness under every Gaussian law and a uniform Gaussian tail bound because \(V\preceq v_{\max}I_d\).

Now let \((h_m,V_m)\to(h,V)\), use the continuous positive-definite square root, and couple \(Z_{V_m}=V_m^{1/2}g\), \(Z_V=V^{1/2}g\) with \(g\sim N(0,I_d)\).  The union of the finitely many decision hyperplanes of \(R(h+Z_V)\) has probability zero because \(V\succ0\).  Off that union the winning action in every stratum is locally constant.  On such a realization, split the temperature domain into a small neighborhood of zero, a compact interior, and a neighborhood of one.  The hard bound controls the first piece, ordinary differentiability controls the compact piece, and the high-temperature bound controls the last piece.  It follows that
\[
 \|G(h_m,Z_{V_m})-G(h,Z_V)\|_\infty\longrightarrow0
 \quad\text{almost surely}.
\]
Thus \(S(h_m,Z_{V_m})\to S(h,Z_V)\) almost surely, and probabilities converge on every continuity set.

Finally, the preceding uniform tail bound makes all individual right \((1-\alpha)\)-quantiles finite uniformly over the compact parameter set.  The least-favorable quantile is their supremum: indeed, \(\inf_{h,V}P\{S(h,Z_V)\le q\}\ge1-\alpha\) holds exactly when \(q\) is at least every individual right quantile.  Therefore \(Q_{1-\alpha}^{\mathrm{LF}}<\infty\), and at that value every individual closed CDF is at least \(1-\alpha\).  This last conclusion uses right quantiles and remains valid at atoms.

%%% FIELD proof-operational-selector-bridge
Each cell count, cell sum, empirical stratum weight, and hence each coordinate of \(\widehat u_n\) is \(\mathcal F_n^{\mathrm{obs}}\)-measurable.  The formulas defining the plug-in frontier use only finite arithmetic, maxima, exponentials, logarithms, and the fixed benchmark.  The endpoint lemmas show that the resulting pointwise formula is a continuous path.  Rational-time measurability and separability of \(C([0,1])\) then show that \(\widehat{\mathfrak F}_n\) is an \(\mathcal F_n^{\mathrm{obs}}\)-measurable random element.  Its displayed construction contains neither \(\Delta_n\), \(u_n\), nor \(E_n\).

For every path \(f\), every index \(s\in[0,1]\), and every real \(x\), \(\sup_tf(t)\le x\) implies \(f(s)\le x\).  Apply this deterministic fact pointwise with \(f=T_n\) and \(s=\widehat t_n\) to obtain
\[
 \{\sup_tT_n(t)\le x\}\subseteq\{T_n(\widehat t_n)\le x\}.
\]
No independence between \(T_n\) and the selector is used.  If \(0<\widehat t_n<1\), the exact regret identity with \(\widehat c_n=\widehat t_n/(1-\widehat t_n)\) gives
\[
 T_n(\widehat t_n)=(1+\widehat c_n)\sqrt n\,\mathscr R_n(\widehat c_n).
\]
Thus the observed frontier selects the index, while the covered target remains the oracle causal regret.

%%% FIELD proof-uniform-path-law
We first record the elementary coupling consequence of bounded-Lipschitz distance.  If \(d_{\mathrm{BL}}(P,Q)\le r\le1\), then for every closed \(A\) the ramp \((1-d(x,A)/\delta)_+\), with \(\delta=\sqrt r\), gives \(P(A)\le Q(A^\delta)+\delta\), and symmetrically.  Approximate the two laws by finite partitions of mesh tending to zero; the preceding inequalities are precisely the cut inequalities for a bipartite transport that leaves at most \(\delta\) mass outside pairs at distance at most \(2\delta\).  Finite max-flow/min-cut and weak compactness of couplings therefore give random vectors \((E,Z)\) with the prescribed laws such that
\[
 P(\|E-Z\|_\infty>2\sqrt r)\le2\sqrt r.                \tag{2}
\]
This derivation also shows that the same coupling can be retained when the vector itself is appended to any mapped process.

Apply (2) with \(r=r_n\), \(E=E_n\), and \(Z=Z_{V_n}\).  Put \(x=2\sqrt{r_n}+\rho_n\); changing the harmless factor two only enlarges the final constant multiplying the frozen \(x_n=\sqrt{r_n}+\rho_n\).  For \(M\ge1\) and \(0<a\le1/4\), let \(\mathcal G(M,a)\) be the event that \(\|Z\|_\infty\le M\), the coupling error is at most \(2\sqrt{r_n}\), and every augmented action-coordinate gap of \(R(h+Z)\) exceeds \(a\).  The hard-endpoint slab bound and the Gaussian coordinate tail give
\[
 P\{\mathcal G(M,a)^c\}
 \le C_0\left[\sqrt{r_n}+\frac{JK^2a}{\sqrt{v_{\min}}}
 +d e^{-M^2/(2v_{\max})}\right]                       \tag{3}
\]
for the numerical constant \(C_0=4\).  The condition \(x_n\le a/8\) ensures that the coupled, locally shifted argument remains in the same hard cell.

Here is the deterministic path estimate used on \(\mathcal G(M,a)\).  Let
\[
 c_0=\frac{a}{8\log\{\exp(1)+M/a\}}.
\]
For \(c\le c_0\), compare both paths with their common hard-cell formula.  The hard-endpoint estimate, \(c\ell_c\le c\operatorname{osc}(Rz)\), and \(e^{-a/(2c_0)}\le\{\exp(1)+M/a\}^{-4}\) show that the total error is at most \(C_1(Ma+x)\).  For \(c\ge c_0\), use
\[
 D_u\ell[r]=\{q(u)-q(u+z)\}^\top r+c^{-1}z^\top C\{q(u+z)\}r,
 \qquad
 D_z\ell[r]=c^{-1}z^\top C\{q(u+z)\}r,
\]
when \(c\le1\), and use the integral covariance representation when \(c\ge1\).  The coordinate derivatives of \(C(q)\) are bounded by \(6/c\).  The mean-value theorem consequently gives
\[
 \sup_{c\ge c_0}|(1+c)\ell_c(u',z')-(1+c)\ell_c(u,z)|
 \le C_2\frac{x(1+M)^2\log\{\exp(1)+M/a\}}{a}.        \tag{4}
\]
The constants \(C_1,C_2\) depend only on \(J,K,\beta\) after summing with the weights \(p_j\); the endpoint at infinity is included by the high-temperature expansion.  Equations (3)--(4), the bound by one in the definition of the bounded-Lipschitz metric, and the coupling give
\[
 d_{\mathrm{BL}}\!\left(\mathcal L\{G(u_n,E_n)\},
                    \mathcal L\{G(h,Z_{V_n})\}\right)
 \le C_{\mathrm{map}}\left[\sqrt{r_n}+\frac{JK^2a}{\sqrt{v_{\min}}}
 +d e^{-M^2/(2v_{\max})}+Ma
 +\frac{x_n(1+M)^2\log\{\exp(1)+M/a\}}a\right],
\]
where, for example, tracing the displayed bounds permits
\[
 C_{\mathrm{map}}=2^{12}JK^2\beta^{-2}p_{\min}^{-1}
 (1+v_{\min}^{-1/2}+v_{\max}+\operatorname{diam}(\mathcal H)).
\]
Taking the stated infimum proves the advertised bound.  Applying the same argument to a bounded one-Lipschitz test of \((G(u_n,E_n),E_n)\) proves the joint assertion.

It remains to check convergence.  Set \(M_n=\max\{1,\sqrt{4v_{\max}\log(1/x_n)}\}\) and \(a_n=x_n^{1/3}\) once \(x_n<1\).  For all sufficiently large \(n\), \(x_n\le a_n/8\); every displayed term then tends to zero, including
\[
 x_n a_n^{-1}(1+M_n)^2\log\{\exp(1)+M_n/a_n\}.
\]
Since \(r_n\to0\) and \(\rho_n\to0\), this proves \(\varepsilon_n\to0\) uniformly over the stated model class.

%%% FIELD proof-certified-calibration
Let \(\delta_n=\alpha-\sqrt{\overline\varepsilon_n}-\gamma>0\).  From the definition of \(B_\gamma\),
\[
 2d\exp\{-B_\gamma^2/(2v_{\max})\}=\delta_n.
\]
For every \(V\preceq v_{\max}I_d\), the scalar Gaussian tail bound and a union bound imply
\[
 P(\|Z_V\|_\infty>B_\gamma^+)\le
 2d\exp\{-(B_\gamma^+)^2/(2v_{\max})\}\le\delta_n,     \tag{5}
\]
because directed rounding certifies \(B_\gamma^+\ge B_\gamma\).  The global Bregman and hard-loss bounds give, for every \(h,z\),
\[
 S(h,z)\le2\|z\|_\infty+\tfrac12\|z\|_\infty^2.      \tag{6}
\]
Indeed \(\operatorname{osc}(Rz_j)\le2\|z\|_\infty\), \(c\ell_c\le\operatorname{osc}(Rz_j)^2/8\), the hard loss is at most the same oscillation, and the quadratic endpoint is at most \(\operatorname{osc}(Rz_j)^2/8\).  Equations (5)--(6) prove that the rational fallback
\[
 q^+=2B_\gamma^++(B_\gamma^+)^2/2
\]
has Gaussian lower CDF at least \(1-\alpha+\sqrt{\overline\varepsilon_n}+\gamma\).  After the probability ledger spends at most \(\gamma\), the required lower bound remains.

On the refined branch, every retained parameter box is an outer box, every discarded covariance box has a directed spectral-disjointness witness, every accepted Gaussian box is contained in the sublevel event after the three deterministic off-grid errors, and every density integral is rounded downward.  Therefore the exact-rational sum \(L_{\mathrm{CDF}}\), after subtracting the three probability-ledger upper bounds, is no larger than the true worst-case CDF.  The verifier's inequalities consequently imply
\[
 \inf_{h\in\mathcal H,V\in\mathcal K_V}
 P\{S(h,Z_V)\le Q^+_{n,\alpha}+\zeta\}
 \ge1-\alpha+\sqrt{\overline\varepsilon_n}.            \tag{7}
\]
This is a one-sided inclusion argument and never inverts a continuous approximation to a CDF, so atoms cause no defect.

The algorithm performs no more than \(w_{\max}\) refinement nodes.  Each node has at most \(K^J\{8+4d+d(d+1)\}\) certified box primitives, the fallback and initialization use at most \(1+K^J\{8+4d+d(d+1)\}\), and radius evaluation uses at most \(4d(m_{\mathsf A}+1)\) interval calls.  Their sum is bounded by the stated \(N_{\mathsf A}\).  Totality and the width guarantee of the interval oracles imply that every call at precision at most \(m_{\mathsf A}\) halts; after the work cap, the fallback just proved is emitted.  Rational range reduction plus Taylor remainders evaluates each elementary function with bit cost polynomial in the input and requested output bit lengths.  Thus the claimed unit-cost count and bit-complexity statement follow, and every accepted certificate has implication (7).

%%% FIELD proof-atom-safe-envelope
Put \(q_n=Q^+_{n,\alpha}+\zeta\) and \(\delta_n=\sqrt{\overline\varepsilon_n}\).  Because \(\overline\varepsilon_n<1\), define the one-sided ramp on path space by
\[
 \varphi_n(f)=
 \begin{cases}
 1,&\sup_tf(t)\le q_n,\\
 1-\{\sup_tf(t)-q_n\}/\delta_n,&q_n<\sup_tf(t)<q_n+\delta_n,\\
 0,&\sup_tf(t)\ge q_n+\delta_n.
 \end{cases}
\]
The supremum map is one-Lipschitz, so \(0\le\varphi_n\le1\) and \(\varphi_n\) is \(\delta_n^{-1}\)-Lipschitz.  The uniform path law and \(\varepsilon_n\le\overline\varepsilon_n\) therefore give, for all sufficiently large \(n\),
\[
 E_{P_n}\varphi_n(T_n)
 \ge E\varphi_n\{G(h,Z_{V_n})\}
      -\varepsilon_n\max(1,\delta_n^{-1})
 \ge E\varphi_n\{G(h,Z_{V_n})\}-\delta_n.             \tag{8}
\]
The certified calibration gives
\[
 P\{S(h,Z_{V_n})\le q_n\}\ge1-\alpha+\delta_n
\]
uniformly because \(h\in\mathcal H\) and \(V_n\in\mathcal K_V\).  Since the Gaussian sublevel indicator is bounded above by \(\varphi_n\), while \(\varphi_n(T_n)\) is bounded above by the indicator of \(\sup_tT_n(t)\le q_n+\delta_n=C_{n,\alpha}\), (8) yields
\[
 \inf_{(P_n)\in\mathcal M}P_n\{T_n(t)\le C_{n,\alpha}
       \text{ for all }t\in[0,1]\}\ge1-\alpha
\]
for all sufficiently large accepted rows.  Taking the requested lower limit proves the first assertion.  Notice that no probability of an interval \((q_n,q_n+\delta_n]\) was bounded; an atom is therefore harmless.

For every observed-data-measurable selector, the deterministic inclusion from the selector bridge sends the simultaneous event into \(\{T_n(\widehat t_n)\le C_{n,\alpha}\}\), without an independence assumption.  The same lower bound follows.  On the interior, the exact regret identity changes this event, and only this event, into
\[
 \{(1+\widehat c_n)\sqrt n\,\mathscr R_n(\widehat c_n)\le C_{n,\alpha}\}.
\]
This proves selection-valid coverage for the causal oracle regret while keeping the plug-in frontier solely on the selection side.

%%% FIELD proof-covariance-inflation-sharpness
In this submodel write \(z\in\mathbb R\) and \(S_0(z)=S(0,z)\).  Benchmark symmetry makes \(S_0\) even.  For \(z>0\) and finite \(c>0\), direct differentiation gives
\[
 \frac{d}{dz}\ell_{c,b}(0,z)=\frac{z}{c}q_{c,b}(z)\{1-q_{c,b}(z)\}>0,
\]
and the quadratic endpoint has derivative \(z/4>0\).  The hard endpoint is zero.  Since every positive \(z\) gives a positive interior or quadratic value, no maximizer can lie at the hard endpoint.  Compactness of \([0,1]\) then shows that \(S_0\) is continuous and strictly increasing in \(|z|\).

Let \(a_\alpha=\Phi^{-1}(1-\alpha/2)>0\).  Monotonicity and continuity give the exact right-quantile formula
\[
 q_{1-\alpha}(V)=S_0(a_\alpha\sqrt V).
\]
Fix an interior \(V\) and put \(z_0=a_\alpha\sqrt V\).  The set of temperatures attaining, or lying within half the positive value \(S_0(z_0)\), is a compact subset of \((0,1]\).  On a neighborhood of \(z_0\), the displayed derivative and the quadratic derivative have a common positive lower bound \(m\) and a finite upper bound \(M\) on that compact temperature set.  Maximizers at nearby \(z\) stay in the same set.  The mean-value theorem for each active coordinate therefore yields
\[
 m|z-z'|\le |S_0(z)-S_0(z')|\le M|z-z'|
\]
for \(z,z'\) sufficiently close to \(z_0\).  Finally,
\[
 \sqrt V-\sqrt{V-s}=\frac{s}{\sqrt V+\sqrt{V-s}},
\]
whose denominator is bounded above and away from zero for small \(s>0\).  Combining the last two displays supplies constants \(c_1,c_2,s_0>0\) with the asserted two-sided linear inequality.  Hence an undershoot of order \(s\) changes the exact critical value by order \(s\), proving that a uniform \(o(s)\) inflation cannot cover this legal submodel.

%%% FIELD statement-no-uniform-modulus
For fixed \(h\in\mathcal H\), define the open hard cells \(D_{\mathbf a}(h)=\{z:\ a_j\text{ is the unique maximizer of }R(h_j+z_j)\text{ for every }j\}\), their hard levels \(L_{\mathbf a}(h)=\sum_jp_j[\max_r(Rh_j)_r-(Rh_j)_{a_j}]\), and \(U_h(z)=\sup_{0<t\le1}G(h,z)(t)\).  Then every atom of \(S(h,Z_V)\) is at one of the finitely many values \(L_{\mathbf a}(h)\), and
\[
 P\{S(h,Z_V)=q\}=\sum_{\mathbf a:L_{\mathbf a}(h)=q}P\{Z_V\in D_{\mathbf a}(h),\ U_h(Z_V)\le q\}.
\]
The frozen assumptions do not imply a uniform atom-removed modulus.  In particular, there are legal fixed inputs with \(J=K=2\), \(p=(1/2,1/2)\), common benchmark \(b=(64/65,1/65)\), \(\beta=1/65\), \(\mathcal H=[1,5/4]^2\), and \(\mathcal K_V=\{I_2\}\), constants \(\eta>0\), and sequences \(h_m\in\mathcal H\), \(q_m\in\mathbb R\), and \(x_m\downarrow0\) such that
\[
 P\{q_m<S(h_m,Z_{I_2})\le q_m+x_m\}\ge\eta.
\]
Consequently no \(A(x)\downarrow0\) having the requested uniform property follows in general, even after the atom at \(q\) is removed.  The one-sided probability certificate of \(\mathsf A\) remains valid, including at all enumerated atoms, but the frozen assumptions alone do not supply a two-sided least-favorable quantile bracket of prescribed value width.

%%% FIELD proof-no-uniform-modulus
Fix \(h\) and a hard cell \(D_{\mathbf a}(h)\).  On that cell the hard endpoint is the constant \(L_{\mathbf a}(h)\), and continuity in temperature gives
\[
 S(h,z)=\max\{L_{\mathbf a}(h),U_h(z)\},\qquad
 \lim_{c\downarrow0}(1+c)\sum_jp_j\ell_{c,b_j}(Rh_j,Rz_j)=L_{\mathbf a}(h).       \tag{9}
\]
Consider a ray \(z=sv\), \(s>0\), with \(v\ne0\).  For every finite \(c>0\), differentiation under the finite sum gives
\[
 \frac{d}{ds}\left[(1+c)\sum_jp_j\ell_{c,b_j}(Rh_j,sRv_j)\right]
 =\frac{s(1+c)}c\sum_jp_j(Rv_j)^\top C\{q_{c,b_j}(R(h_j+sv_j))\}Rv_j>0.          \tag{10}
\]
The inequality is strict because some block \(v_j\ne0\), \(Rv_j\) is then nonconstant, every Gibbs coordinate is positive, and a categorical variance vanishes only for a constant vector.  The quadratic endpoint is also strictly increasing along the ray.  The intersection of an open polyhedral cell with a ray is an interval.  If \(s_1<s_2\) lie in this interval and \(U_h(s_1v)>L_{\mathbf a}(h)\), compactness of the temperature compactification and (9) imply that the supremum at \(s_1\) is attained at some \(t>0\); (10) then gives \(U_h(s_2v)>U_h(s_1v)\).  Therefore a positive-length flat radial interval wholly inside the cell can only have value \(L_{\mathbf a}(h)\), and, for every \(q\ne L_{\mathbf a}(h)\), each radial section of
\[
 D_{\mathbf a}(h)\cap\{z:U_h(z)=q\}
\]
has one-dimensional Lebesgue measure zero.

Polar Fubini now makes that level set Lebesgue null.  The finitely many decision boundaries are affine hyperplanes and are Gaussian null because \(V\succ0\).  Since the Gaussian has a Lebesgue density, cells with \(L_{\mathbf a}(h)\ne q\) contribute no mass to \(\{S=q\}\): if \(L_{\mathbf a}>q\) the event is empty, and if \(L_{\mathbf a}<q\) it is contained in the null set \(\{U_h=q\}\).  On a cell with \(L_{\mathbf a}=q\), (9) says exactly that \(S=q\) iff \(U_h\le q\).  Summation over the finitely many cells proves the atom formula.

It remains to show that these atoms prevent a uniform atom-removed modulus for some legal fixed inputs.  Take the inputs in the statement and write \(r=b_1/b_0=1/64\).  Both coordinates of \(h=(h_1,h_2)\in[1,5/4]^2\) have action \(1\) as the true maximizer.  At the point \(z^{(1)}=(-2h_1,0)\), only stratum one is assigned to the wrong hard action, and the hard level is \(L_1(h)=h_1/2\).  At \(z^{(2)}=(0,-2h_2)\), only stratum two is wrong, and the hard level is \(L_2(h)=h_2/2\).

We verify that each point is interior to a positive-mass flat region; this is the required collision argument over all temperatures.  For one wrong stratum put \(x=h/c\).  Direct substitution into the binary log-sum-exp formula gives
\[
 \frac{\ell_{c,b}(h,-2h)}h
 =1+\frac1x\log\frac{r+e^{-x}}{1+re^{-x}}
   -\frac{2re^{-x}}{1+re^{-x}}.                         \tag{11}
\]
For \(0<h\le-\log r\), multiplying (11) by \(1+h/x\), separating the cases \(x\le-\log r\) and \(x\ge-\log r\), and applying \(y/(1+y)\le\log(1+y)\le y\) to the two logarithms gives
\[
 (1+c)\ell_{c,b}(h,-2h)<h\qquad(c>0).                  \tag{12}
\]
Here \(-\log r=\log64>5/4\), so (12) is uniform over the displayed range of \(h\) on every compact temperature interval.

The strict inequality must persist in a neighborhood uniformly also at the two ends of the temperature scale.  Near zero, direct log-sum-exp asymptotics, uniform for \(h\in[1,5/4]\), a wrong estimated contrast bounded above by a negative constant, and a correct estimated contrast bounded below by a positive constant, give respectively
\[
 (1+c)\ell_{c,b}(h,z)=h+c\{h+\log(b_1/b_0)\}
   +c^2\log(b_1/b_0)+O\{(1+c)(c+|z|)e^{-\kappa/c}\},
 \qquad
 (1+c)\ell_{c,b}(h,z)=O\{(1+c)(c+|z|)e^{-\kappa/c}\}. \tag{13}
\]
The first coefficient is at most \(5/4-\log64<0\).  Thus a linear-in-\(c\) wrong-action deficit dominates the exponentially small correct-action loss.  On an intermediate compact interval, (12), compactness in \(h\), and continuity provide a strict uniform margin.  At infinity, the high-temperature expansion gives, at the two central points,
\[
 \lim_{c\to\infty}(1+c)\ell_{c,b}(h,-2h)=2b_0b_1h^2<h,
\]
uniformly for \(h\in[1,5/4]\); its \(O(c^{-1})\) remainder and continuity give the same inequality on a neighborhood.  Combining the three temperature ranges produces open boxes \(B_1(h)\ni z^{(1)}\), \(B_2(h)\ni z^{(2)}\), with radii bounded below uniformly in \(h\), on which \(S=L_1(h)\) and \(S=L_2(h)\), respectively.  The standard Gaussian density is bounded below on the resulting compact union, so
\[
 \inf_{h\in[1,5/4]^2}P\{Z_{I_2}\in B_k(h)\}\ge\eta>0,qquad k=1,2.             \tag{14}
\]

Take \(h_m=(1,1+\delta_m)\) with rational \(\delta_m\downarrow0\), \(q_m=1/2\), and \(x_m=\delta_m/2\).  The second atom lies at \(q_m+x_m\), while the first lies at the excluded left endpoint \(q_m\).  Equation (14) yields
\[
 P\{q_m<S(h_m,Z_{I_2})\le q_m+x_m\}\ge\eta.
\]
These fixed inputs are compatible with the primitive model: take \(p_j=e_{ja}=1/2\), independent two-point potential-outcome noises of variance \(1/8\), means \(\mu_{n,j0}=0\), \(\mu_{n,j1}=h_j/\sqrt n\), and consistency.  Then bounded outcomes hold, the local remainder is zero, and the two contrast scores have covariance \(I_2\).  Hence the witness is legal.  The lower bound contradicts every proposed \(A(x)\downarrow0\).  Finally, the certified-calibration theorem uses only a lower CDF bracket and its analytic fallback, so its one-sided conclusion survives these atoms; only the requested uniform route to a two-sided value-width bracket fails.
